% ************************************************************
% INTRODUCTION
% ************************************************************
\section{Introduction}

WASH disruption in Gaza is a public-health systems problem involving water access, sanitation, environmental exposure, displacement, health-service access, and operational continuity. Public humanitarian reporting in April and May 2026 documented concurrent displacement-site environmental hazards, dependence of WASH services on fragile generator systems, sewage and wastewater exposure, and constrained health-service capacity \citep{ocha_2026_apr17,ocha_2026_may15}. In settings where primary data collection is infeasible or unsafe, public humanitarian reporting can support rapid prioritization if extraction rules, scoring assumptions, and evidence limits are explicit.

This study develops a first-pass secondary-data index intended for refinement through verified WHO, WASH Cluster, OCHA, and HDX data layers. The contribution is a reproducible prioritization framework rather than a causal estimate of disease burden. The index is designed to identify where WASH service stress, environmental-health hazards, health-system access constraints, population vulnerability, and operational constraints converge in publicly available data.
% ************************************************************
% END INTRODUCTION
% ************************************************************
